// Caustics — scatter-write underwater light simulation
//
// Uses scatter writes (@OUT[x, y] += value) to simulate underwater
// caustics. A noise-based displacement field bends "light rays" and
// scatters energy to displaced positions. Where rays converge, energy
// accumulates into bright caustic patterns.
//
// This is a scatter operation: each pixel computes WHERE its energy
// lands, not what lands ON it. The += accumulation creates the
// characteristic bright convergence lines of real caustics.
//
// No input image required — generates a procedural caustic pattern.

f$frequency = 3.0;      // noise frequency (higher = finer pattern)
f$amplitude = 12.0;     // displacement strength in pixels
f$energy = 0.08;        // base light energy per pixel
f$time = 0.0;           // animation offset (connect frame index)
f$warmth = 0.3;         // warm tint amount (0 = white, 1 = aqua-gold)

// ── Initialize output to black ──────────────────────────────────────
@OUT[ix, iy] = vec3(0.0);

// ── Compute noise-based displacement ────────────────────────────────
// Two layers of simplex noise create a 2D vector field that bends
// the light rays. The offset between noise samples creates curl-like
// swirling patterns typical of water caustics.
float nx = u * $frequency + $time * 0.1;
float ny = v * $frequency + $time * 0.07;

float dx = simplex(nx, ny) * $amplitude;
float dy = simplex(nx + 5.3, ny + 7.1) * $amplitude;

// Second octave for fine detail
dx += simplex(nx * 2.0, ny * 2.0) * $amplitude * 0.5;
dy += simplex(nx * 2.0 + 5.3, ny * 2.0 + 7.1) * $amplitude * 0.5;

// ── Scatter energy to displaced position ────────────────────────────
// The convergence of displaced rays creates bright caustic lines.
int tx = int(ix + dx);
int ty = int(iy + dy);

// Tint: blend between white and a warm aqua-gold
vec3 warm = vec3(0.9, 1.0, 0.8);
vec3 cool = vec3(0.8, 0.95, 1.0);
vec3 tint = lerp(vec3(1.0), lerp(cool, warm, luma(vec3(u))), $warmth);

@OUT[tx, ty] += tint * $energy;
